\documentclass{jsarticle}
\usepackage[dvipdfmx,bookmarks=true]{hyperref}
\usepackage{pxjahyper}
\usepackage{pdfpages}
\usepackage{ulinej}
\usepackage{url}
\hypersetup{
 setpagesize=false,
 bookmarksnumbered=true,%
 colorlinks=true,%
 linkcolor=blue,
 citecolor=red,
}
\usepackage{ascmac}

\title{不明点に対する回答書}
\author{\empty}
\date{\empty}

\begin{document}
\maketitle
\noindent \textbf{論文番号：P45} \\\
\noindent \textbf{標題：家庭環境知識に基づく行動前提条件の検証を伴うLLMエージェントの逐次的行動計画生成}  \\

査読いただき，ありがとうございます．本回答では，保存済み結果の再集計，公開実装の仕様確認，新しいログ付き実験を区別して報告します．既存結果は状態変化312件・配置103件に対する18条件であり，追加比較は4タスク・2モデル・10条件・3反復の240エピソードです．根拠資料のパスはリポジトリルートからの相対パスで示します．論文中の参照先は節名・表題で示し，回答書内の表番号とは区別します．



\newpage

\section*{査読者A様の査読結果への回答}

\subsection*{査読者A様のコメント}
\begin{screen}
論文としての体裁や構成はしっかりしており、読みやすいです。
実験も十分にされております。
しかしながら、新規性に問題があります。

この論文に新規性があるとすれば、
「次のアクションの事前条件が満たされていない場合は、別のアクションをLLMに再生成させる。」
という点につきるかと考えられます。
論文中でも、この点を意識して、比較実験をされています。

しかしながら、一般のオンラインプランニングでは、
各アクションの事前条件が満たされているかどうかプラン実行中にチェックして、
満たされていない場合は現在の状態の記録（信念）を更新した上で、
再プランニングすることはよく行う手段であり、
ここに新規性があるとはいえないです。
特にLLMを用いる上で特別な工夫点が必要というわけでもなさそうです。

従いまして、内容的には一つの事例研究としては面白いのですが、
通常の原著論文としては、この新規性の欠如は致命的で、容易に改善できないと思われるので、
今回は不採録判定とさせていただきます。

既に実験環境やプログラムは整っており、新たな実験はやりやすいかと思います。

論文のストーリーを検討しなおし、技術面での新規性を明確化し、
必要に応じて新しい技術要素を追加したり、
応用面で得られた新しい重要な知見を入れたりして、
それに合わせて評価実験を行い、論文をまとめなおし、
再投稿することをお勧めいたします。
\end{screen}

\subsection*{査読者A様のコメントへの回答}
論文の位置付けと評価の方向性について，具体的なご指摘をいただき，ありがとうございます．ご指摘の通り，行動の前提条件を確認し，現在状態に応じて計画を修正することは，オンラインプランニングにおける既存の考え方です．この点を踏まえ，改訂では，既存の実行監視・計画修復を基礎として，家庭環境知識を用いるLLM逐次計画の具体的な構成と，各処理段階の働きを明らかにする研究として論文を再構成しました．特に，ご提案いただいた「応用面で得られた新しい重要な知見」を具体化するため，構成要素を切り替える追加比較と，実行ログに基づく過程別の分析を加えました．

まず，既存研究との関係を明確にするため，第2章にPLANEX1による記号的実行監視とLPGによる計画修復を追加しました．また，Huang・Raman・Linらと本研究を，入力情報，検証・照合時点，判定主体，修正方法，評価条件の共通軸で整理しました．Linらの反復デコードと実行エラーに基づく再計画についても，それぞれの修正の契機を明示しました．これにより，本研究が継承する一般原理と，本研究で具体的に評価する構成との対応を示しました．

次に，技術的構成と評価対象を具体化しました．本研究では，同じ指示でも初期状態に応じて行動を変更する家庭内タスクを扱います．例えば，``Turn on all lightswitches''に対して，同名のスイッチをIDで識別し，現在のON／OFF状態に応じて操作対象を選びます．この判断を支えるため，タスクに関係する物体の状態・ID・支持／包含関係などを抽出して自然言語化し，候補行動に対応する前提条件とともにLLMへ提示します．さらに，前提条件判定，未充足条件の説明，1回の再生成，終了判定，実行結果を対応付けて追跡できるよう，実装仕様を記載しました．状態変化・配置の計618インスタンスについても，対象範囲・分割・参照手順を明示しました．これらの実装とタスク集合を，LLMを組み込んだ逐次計画の振る舞いを検証する基盤として位置付けています．

追加評価では，「前提条件の判定をLLMに担わせる効果」「実行前の検証と失敗後の修復の違い」「前提条件の明示と一般的な再考の違い」を調べる対照条件を設けました．状態変化\texttt{test\_task1, test\_task6}と配置\texttt{test\_task1, test\_task65}の4タスクを，初期充足，未達対象の追跡，複数物体の配置，容器操作という判断を調べるために選び，2モデル・10条件・各3反復の計240エピソードを実行しました．表\ref{tab:reply_controls}に成功数を示します．モデルIDは\texttt{gpt-4o-mini-2024-07-18}と\texttt{gpt-4o-2024-08-06}，temperatureは0です．全240件で，API呼出前に参照との記号状態一致，全ノード位置差0.001以下，四元数距離0.0001以下を照合しました．初期充足の1タスクは全条件で成功し，各分子に3件ずつ含まれます．参照列の再生が失敗した配置1件も集計対象とし，参照成功3件だけの集計を別に保存しています．

\begin{table}[htbp]
\centering\small
\caption{追加比較の成功数／12（4タスクの各3反復）}\label{tab:reply_controls}
\begin{tabular}{lrr}
\hline
条件 & GPT-4o mini & GPT-4o \\
\hline
C0：検証なし & 8/12 & 9/12 \\
C1：LLMによる実行前検証 & 9/12 & 9/12 \\
C2：抽出知識の記号検証 & 4/12 & 10/12 \\
C3：VH失敗後の修復 & 5/12 & 9/12 \\
C1-budget：条件レビューと毎回再生成 & 4/12 & 11/12 \\
C4-budget：一般的再考と毎回再生成 & 6/12 & 10/12 \\
C5：完全グラフの記号検証 & 5/12 & 9/12 \\
R-random：C1の事例をランダム選択 & 8/12 & 9/12 \\
R-fixed：C1の事例を固定 & 7/12 & 9/12 \\
R-none：C1の事例を除去 & 3/12 & 3/12 \\
\hline
\end{tabular}
\end{table}

この追加評価から，以下の三点を診断的知見として整理しました．

第一に，検証方式と再生成機会の効果は，モデルと条件に応じて評価する必要があります．LLM検証C1と検証なしC0の成功数の差はminiで1件，4oで0件でした．同じ抽出知識を用いる記号検証C2とC1の比較ではminiで5件減，4oで1件増となりました．また，毎行動決定の生成・レビュー・説明・再生成の4要求と出力上限合計1024トークンをそろえたC1-budgetとC4-budgetでは，前提条件を明示する側の成功数はminiで2件少なく，4oで1件多くなりました．これらを，判定方法と再考の機会を分けて検討する根拠として示しました．この対照でそろえた計算量の範囲は要求数と出力上限であり，実入力長・実出力長・エピソード全体の要求数は異なります．C5についても，完全グラフの参照に加え，既知の空状態集合の扱いがC2と異なる条件として説明しました．

第二に，実行前検証と失敗後修復の違いを捉えるには，最終的な成功率と途中の失敗実行数を併せて評価することが有用です．C3はVH失敗後の最新状態とエラーに基づいて1回修復する独自対照です．GPT-4oではC1とC3の成功数はともに9/12でしたが，途中を含む失敗実行数はC1で3回，C3で9回でした．この観測から，同じ成功数に至る場合にも実行過程に違いがあることを示しました．

第三に，LLMによる条件判定，終了判定，シミュレータでの実行可能性を分けて診断することで，改善すべき処理段階を具体化できます．C1のLLM判定と同じ候補・抽出辞書上のJSON条件との一致数は，miniで34/43，4oで45/54でした．全10条件の終了判定では，miniは496回中26回で未達のままEndを返しました．また，完全グラフ上のJSON条件がtrueだった実行でも，miniは376回中28回，4oは447回中37回でVHが失敗を返しました．各段階の入力・応答と実行前後の状態を保存し，条件判定の誤り，未達対象を残す終了，辞書の条件を満たした行動の実行失敗を区別して説明しました．具体的な判定精度とケースはB-4-3への回答にも示します．

以上を踏まえ，第1章の貢献範囲，第2章の表「関連研究との評価対象の比較」，第5章「構成要素の比較」，第6章および英文要旨を改訂しました．本研究の貢献を，初期状態に応じた判断を扱うタスク集合，家庭環境知識とLLM逐次計画を接続する具体的実装，その構成要素の働きと失敗過程を明らかにする知見として提示します．追加比較の結論は，目的選択した4タスク（指示名3種類）の各3反復で観測した範囲に位置付け，全件反復評価，厳密な計算量一致対照，実環境の安全性評価を今後の評価項目として明示します．

追加比較の根拠は\path{Paper/vh_step4/metrics.csv}，\path{Paper/vh_step4/paired_results.csv}，\path{Paper/vh_step4/repeat_metrics.csv}および同ディレクトリの\path{config.json}・\path{runs/}です．判定・終了・実行の過程別集計は，同ディレクトリの\path{precondition_metrics.csv}，\path{termination_labels.csv}，\path{execution_labels.csv}に示します．

\newpage

\section*{査読者B様の不明点への回答}

\subsection*{不明点B-1:タスク設計が不明瞭}
\subsubsection*{不明点B-1-1}
\begin{screen}
タスク記述（文）はどのような粒度での設定を目指したのか、そのポリシーが不明確です。``Turn on all lightswitches''が今回設定したタスク記述だと思われるのですが、``Turn on all lightswitches in the kitchen'', ``Turn on all 4 lightswtiches''などではダメなのか、設定のポリシーがよくわかりません。「人間が指示を出すときに、暗黙の裡に省略を行ったようなタスク指示文」を想定しているのではないか、と漠然と読めるのですが、その省略がどの程度なのか、方針を示してほしいと考えます。例えば、人間だったらほぼ間違えないレベルなのか、それとも同居人のようにある程度コンテクストを共有していないとわからない指示文があるのか、といった方針が気になります。
\end{screen}

\subsubsection*{不明点B-1-1への回答}
指示の省略範囲を明確にするため，第3章「指示文の範囲と粒度」に規約を明記しました．指示は状態変化の``Turn on/off all [対象クラス]''と配置の``Put all [対象クラス] in/on the [配置先クラス]''という定型で，一つの対象クラスと，配置では一つの配置先を扱います．``all''は初期部屋を含む当該シーン内の対象クラス全体を指します．ID・個数・現在状態・所属部屋は環境知識から求めます．今回の指示集合はこの定型を対象とし，``in the kitchen''や``all 4 lightswitches''は対象範囲や個数情報を変更する別条件として位置付けます．公開コードでは，この規約を家屋全体から名詞に対応する物体を抽出する処理として実装しています．根拠は\path{dataset/state_change_task/test.json}，\path{dataset/placement_task/test.json}です．代表ケースで必要な判断を説明します（B-4-2）．

\subsubsection*{不明点B-1-2}
\begin{screen}
タスクの複雑さに関するポリシーが不明確です。目標遂行のために、途中にある障害物を移動させることや、一度後戻りするような状態遷移が必要になることがないようにタスクを設計したと思われますが、そのようなポリシーを明確に記載してください。どの程度の推論の複雑さまでを求めるタスクなのか、整理してください。
\end{screen}

\subsubsection*{不明点B-1-2への回答}
第3章「推論の範囲とゴール以外の制約」に，同名物体の識別，現在状態に応じた対象選択，支持・包含関係，把持・近接・容器開放の判断を対象とすることを記載しました．第5章の表「テスト集合の規模と参照行動列長」では，ゴール要素数を状態変化2--5，配置2--6，参照列長の範囲をそれぞれ0--10，0--25と示しています．複雑さは，現存する資料で確認できるこれらの属性に基づいて説明します．障害物移動や達成済みゴールを崩す必要性の分類は，注釈・構築記録の整備を要する項目として扱い，人間にとっての難しさと区別しました．根拠は\path{Paper/analysis/task_inventory.csv}と\path{Paper/analysis/dataset_summary.csv}です．
% 本不明点については，不明点B-2-1について，加筆・修正すれば，明確になるように思います．


\subsubsection*{不明点B-1-3}
\begin{screen}
タスク内でどこまでの情報を与え、どこまでが必要になるタスクなのか、そのポリシーが不明確です。Listing 6のタスクについて、``Turn on all lightswitches in the house''のように、複数の場所を移動することを想定するタスク「ではない」と、この指示を出されたエージェントはどこから判断できるのでしょうか？もし、kitchenにいることが特に意味はなく、単にid 173などのスイッチへ移動してONにするだけ、というタスクだとしたら、initial roomの情報は必要なのでしょうか？ id 173などのスイッチの「場所」は、kitchenの中にあると与える前提でしょうか？そもそも「場所」は考えない、という設定でしょうか？ どこまでの情報を入力に入れ、どこからは与えない設定なのか、を明確にしてください。
\end{screen}


\subsubsection*{不明点B-1-3への回答}
第4章「実環境知識の入力例」に，構成要素比較C1・GPT-4o・反復0のtask6の環境入力全文を掲載しました．ID 71はbathroom 11，173はbedroom 73，261はkitchen 205，427はlivingroom 335に属し，人物の初期位置も記載しています．これはこの試行で取得した状態です．
初期部屋はキャラクター配置に用い，対象範囲の限定には用いません．家屋内の複数部屋をまたぐ移動を許し，計画器は対象名・ID付きのWALKを選択します．細かな経路・前進・回転はVH側が処理します．\texttt{scene}と\texttt{initial\_room}のフィールドや位置座標をそのままLLMに渡すのではなく，現在グラフから抽出した所属部屋・近接・左右の保持関係を文として渡します．\texttt{goal\_states}と評価対象の参照行動列は計画プロンプトに渡しませんが，参照列長は試行上限に利用します．第3章の表「データの入力先と用途」と第4章「家庭環境知識の取得」に対応を整理しました．具体的な追加試行の入力は\path{Paper/vh_step4/runs/}の\path{events.jsonl}で確認できます．

\subsubsection*{不明点B-1-4}
\begin{screen}
タスクの成功を判断する際、goal state以外の状態をどのように考えるかが不明確です。Listing 7のタスクは、冷蔵庫を「閉める」ことはgoal state的には必要なさそうです。このタスクの成功判定には、冷蔵庫が閉まっていることは含まれるのでしょうか？ 含まれないとしたら、正解例のaction scriptはどこまでのactionが載っているべきだと考えられているのでしょうか？ あるいは、タスク終了時に部屋の中がめちゃくちゃに散らかっている、みたいな状況については成功と判断するのでしょうか？

また、このaction例について、人間ならおそらく1個目のplumを入れた後に冷蔵庫を閉め、2個目のplumを持ってきて再度開けるようにも思います。アクション途中で冷蔵庫が異常状態にならないことを要求するならば、そのようなアクションが必要になりそうですが、今回のタスクはそのような制約は考えない、ということで良いでしょうか？明確化してください。
\end{screen}

\subsubsection*{不明点B-1-4への回答}
第3章の制約と第5章の配置Listing・評価指標を修正しました．配置\texttt{test\_task65}のゴールは，plum 53・54がfridge 103のINSIDEにあることとして定義します．公開JSONの参照列に合わせ，掲載例を末尾のCLOSEを除いた9行動に訂正しました．採点対象は指定されたゴールの充足とし，冷蔵庫のCLOSED，途中で閉め直すこと，無関係な物体の原状回復，開放時間は採点対象外として整理しました．途中の操作にはVHの実行条件を適用し，生活上の安全性・快適性は今後の評価項目として位置付けます．

通常終了時の採点は，状態変化では指定IDの状態配列の一致，配置では指定ID間の辺の存在を確認します．ただし実行失敗・試行上限到達は通常採点より優先され，参照列長0の初回Endではグラフ採点を省略する特例があります．この順序も第4章の擬似コードと第5章に明記しました．保存されたGPT-4o・Multiの同タスクの10行動列を追加再生したところ，全行動を再生でき，最終ゴールも充足しました（\path{Paper/vh_saved_replay/Report.md}）．

\subsection*{不明点B-2:タスクセット構築が不明瞭}
\subsubsection*{不明点B-2-1}
\begin{screen}
どの程度異なる複数のタスクを作成したのかが不明確です。スイッチのON/OFFの組み合わせを変えただけなのか、部屋や状況を変えたのか、タスクの「文言・指示の出し方」を変えたのか、変えたとしたらどのように変えたのか、などを明確化してください。全体として、どの程度異なるタスクが、どの程度の範囲でカバーされているのか、理解できることが必要です。
\end{screen}

\subsubsection*{不明点B-2-1への回答}
第5章「評価用データセット」に，対象・配置先の種類数，件数，分割，指示と初期状態の対応を記載しました．全618件の組合せは公開データと\path{Paper/analysis/task_inventory.csv}に示しました．
第5章「データセット構築」に，状態変化対象7種類，配置対象12種類，配置先9種類の全種類と件数を列挙しました．状態変化の指示はON/OFFです．1インスタンスは1対象クラスを扱います．同名インスタンスではシーン，初期部屋，指定ノードの状態配列等を変えます．\texttt{initial\_states}は状態配列の上書きであり，物体位置を任意に変更するフィールドではありません．定型文の対象・操作・配置先を変えます．

全618インスタンスのうち状態変化はテスト312・事例用152，配置はテスト103・事例用51です．シーン1--7を含み，観測した組合せを\path{Paper/analysis/dataset_coverage.csv}と\path{Paper/analysis/dataset_dimension_counts.csv}で開示しました．構築過程の説明は現存資料で確認できる範囲に基づき，物体の除外理由，不具合回避，元の作成・検証手順は未確認項目として明示します．追加試行ではcupcakeの参照列にVH失敗がありました．
% 本不明点については，状態変化が可能な複数種類のオブジェクトや位置を変更可能な複数種類のオブジェクトについて，複数種類のオブジェクトが複数個同時に状態や位置が変化することも想定しているかということが明記されていないことが指摘されているように思います．3節「タスクの設定とデータセットの要件」と5.1.3項「データセット構築」に，一つのタスクにおいて対象とするオブジェクトの種類についての制約も追記すると良いです．また，状態変化が可能なオブジェクトの中で，VirtualHomeの不具合により，エージェントが対象オブジェクトにアクションを実行しても，状態が変化しないものを除いて，タスクを設計していたように思います．この点についても，明記した方が良いと思います．具体的に，今回対象としたVirtualHomeにおけるシーンにおいて，どのようなオブジェクトが状態変化が可能で，その中で，どのオブジェクトを対象としたかを5.1.3項「データセット構築」に列挙すると良いと思います（数種類のため例だけでなく，全て記載した方が明確だと思います）．タスクの「文言・指示の出し方」については，記述方針などがあれば記載すると良いです．


\subsubsection*{不明点B-2-2}
\begin{screen}
同一の名前があるようですが、どの程度類似したものが同じ名前で存在するのか、どの程度近いタスクがどれくらい近い名前になっているのか、がわかりません。これは、提案手法で「訓練データセットから、類似度の高いタスク見本を取り出す」ことがどの程度「当たる」と一般的に期待されるのか、を考える重要な要素になります。
\end{screen}

\subsubsection*{不明点B-2-2への回答}
第3章で指示文を「タスク名」，指示とシーン・初期条件等の組を「インスタンス」と定義しました．テスト／事例用の異なるタスク名数は状態変化6／6，配置38／19で，集合間の同名重複はともに0です．一方で同一集合内には，同じ指示に異なるシーンや初期条件を対応させたインスタンスがあります．シーン・操作構造は両集合で共有します．「訓練用」はパラメータ学習用ではなくプロンプト事例の候補集合です．根拠は\path{Paper/analysis/split_overlap.csv}と\path{Paper/analysis/task_inventory.csv}です．

第4章ではタスク名をMiniLMで埋め込む検索と，同名候補から最長列を選ぶ規則を記述しました．たとえば再構成検索で``Turn off all tablelamps''に``Turn on all tablelamps''が類似度0.910で選ばれています．具体的な検索統計と追加実行比較はB-4-5およびコメントB-1への回答に示します．


\subsection*{不明点B-3:提案手法の詳細}
\subsubsection*{不明点B-3-1}
\begin{screen}
家庭環境知識の取得、でどのような自然言語の表現が得られるのか、が不明確です。例えば物体の情報について、「ほかの物体や部屋との関係」の情報がある、と記載されていますが、例えばListing 6のタスクでは「173のスイッチはキッチンの壁にあり、ロボットと同じ床面に属する」のような記載がListing 2のプロンプトに入るのでしょうか？どこまでの情報がプロンプトに与えられるのか、入れるか入れないかの基準は何か、説明してください。Listing 6,7のタスクでの実例を記載してもらえると、理解が容易であるとも考えます。
\end{screen}
\subsubsection*{不明点B-3-1への回答}
第4章「実環境知識の入力例」に，task6とtask65の実際の環境文全文と，行動による状態・関係の変化を追記しました．コードの振る舞いを説明する合成例とは別に，構成要素比較の特定試行へ対応付けています．
第4章「家庭環境知識の取得」に抽出・自然言語化の処理と実入力例を記載し，コード全文を公開リポジトリに置きました．タスク文のNOUNと単数形候補にクラス名が一致するノードを家屋全体から選び，支持物・容器を1段階追加します．名詞抽出はコードで行います．状態，ON／INSIDE，所属部屋，エージェントの近接・保持を文章にしますが，壁面・床面の任意の関係や位置座標は入力しません．INSIDEは部屋と容器で扱いを分け，ONとINSIDEが両方ある場合はONを優先します．状態配列の先頭要素だけを文にし，所属部屋を得られない物体では基本説明文を出力しません．

公開資料の9000番台IDの例は変換を説明する合成入力です．根拠は\path{Paper/listings/knowledge_example_inputs.json}と\path{Paper/listings/knowledge_serialization.py}です．これとは別に，追加240件では実際の完全グラフ，抽出辞書，送信した自然言語を各runに保存しました．
% Listing 2の説明部分に，Listing2の赤字で示している箇所に具体的にどのようなものが与えられるかの具体例や，取得した家庭環境知識をどのように加工して，プロンプトに含めているかの説明を追記すると良いです．

\subsubsection*{不明点B-3-2}
\begin{screen}
アクションの前提条件の与え方が不明確です。4.6節の最後にある、VHのDocumentationやGitHubの情報を用いた、ということではないか思われますが、明確に記載してください。また、どのようにその情報を使ったのか、機械的に表現が取り出せるのか、事前にある程度整形する必要があったのか、といった具体的な手順も明確にしてください。
\end{screen}

\subsubsection*{不明点B-3-2への回答}
第4章「実装」の「公式資料と前提条件辞書の対応」に，VH v2.3の公式Actions文書（\url{http://virtual-home.org/documentation/master/kb/actions.html}）と公式リポジトリの\texttt{resources}（\url{https://github.com/xavierpuigf/virtualhome/tree/master/virtualhome/resources}）を，前提条件の説明と条件文の照合に利用したことを明記しました．Actions文書では8操作の引数・前提条件・実行後の変化を，\texttt{resources}では物体属性と取り得る状態を確認しました．確認日は2026年9月8日，リポジトリの固定版は\texttt{58970fd80951c2eaa1af713e0917d1a105353ad8}です．

第4章の表「8操作の前提条件辞書と公式資料の対応」に，採用した近接・保持・ON／OFF・OPEN／CLOSEDの条件と各資料の対応を示しました．例えばSWITCHONの近接・OFFはActions文書のSwitchOnに，PUTINの保持・配置先への近接・配置先がCLOSEDでないという3条件はPutInに対応します．GRABの非保持条件は，公開Pythonグラフ実行器の\texttt{GrabExecutor.check\_grabbable}にある重複把持の拒否とも照合しました．同コードは実行履歴の把持記録を検査し，本辞書は左右の保持関係を条件に用います．また，公式文書に記載された物体属性・到達可能性・空いた手などの条件も表に併記しました．

実装では，操作ごとの自然言語条件を\path{precondition.json}に固定して保持します．候補の操作名で条件文を選び，対象名とIDをObject／Object1／Object2へ代入し，現在の環境グラフから抽出した知識とともにLLMへ渡します．条件文の選択と代入はコードで行い，LLM応答に部分文字列Yesが含まれるかで受理を決めます．棄却時には別要求で未充足条件を尋ね，1回再生成した行動を再検証せず実行します．JSON条件の記号評価は追加診断のC2／C5およびラベル作成に導入したものです．資料と各条件の対応は\path{Paper/revision-completion/Precondition-Sources.md}に記録しました．
% 4.6節の最後に，アクションの前提条件をVHのドキュメントやGitHubのソースコードなどから，どのような手順で与えたか説明を追記すると良いです．

\subsubsection*{不明点B-3-3}
\begin{screen}
5.2節第2パラグラフ、プロンプトの与え方を分ける方式について、具体的にどのように分けるのか、説明が不足しています。手法が再現できる程度に記述してください。
\end{screen}

\subsubsection*{不明点B-3-3への回答}
第4章「アクションステップ生成」「呼び出しごとの入力と履歴」にroleと順序を記載しました．空白・改行を含む全文は\path{Paper/listings/prompt_spec.py}で公開しました．行動生成のMultiは，system 1件に加え，(1)冒頭の生成指示，(2)許可操作と書式，(3)選択した事例，(4)現在の環境知識，(5)タスクと成功実行履歴，のuser 5件を1回のAPI要求で送ります．Singleはこの5要素をuser 1件に連結しますが，終了・生成の文面とsystem末尾の空白にも差があります．

終了判定，前提条件判定，違反条件説明，再生成は各々別のメッセージ構成であり，これらも\path{Paper/listings/prompt_spec.py}に全文を保存しました．次周回の履歴に加えるのは実行成功した行動だけです．公開Notebookの仕様と実行記録の対応は，第5章「評価間の条件差と記録範囲」に示しました．
% 4.3節ではListing 2の説明はしていますが，プロンプトの与え方を分ける方式については，説明がされていないため，この説明を追加すれば，本不明点の回答になると思います．

\subsection*{不明点B-4:評価の詳細}

\subsubsection*{不明点B-4-1}
\begin{screen}
想定される限界値との差が明らかになるような情報を提供してください。
表1,2のAverage Stepsは、想定される（タスクセットで正解と定義された）アクション列長とどの程度違うのか、理解できる情報が求められます。何らかの統計量などで示してください。なお、成功した時と失敗したときなどに分けて議論できれば、提案手法の振る舞いを理解するために役立つと考えられます。
\end{screen}
\subsubsection*{不明点B-4-1への回答}
第5章「基本性能評価」に，全18条件の成功・失敗別の件数，平均行動列長と同じ集合の参照平均，検証あり／なしの共通成功集合での平均列長を掲載しました．分布・失敗類型・同一IDの差と比などの詳細集計は公開資料にまとめ，本文では結果の解釈に必要な集計値を示しました．
全18条件の保存結果を再集計し，本文のSR・3失敗率・Average Stepsの計90値が小数第3位まで一致することを確認しました．成功・失敗別の長さを\path{Paper/analysis/lengths_by_outcome.csv}に追加しています．表\ref{tab:reply_lengths}はGPT-4o・Multi Prompts・検証ありの例で，参照平均も各行と同じタスク集合から計算しました．

\begin{table}[htbp]
\centering\small
\caption{既存GPT-4o・Multiの保存行動列長と参照列長}\label{tab:reply_lengths}
\begin{tabular}{llrrr}
\hline
タスク & 結果 & 件数 & 保存平均長 & 参照平均長 \\
\hline
状態変化 & 成功 & 279 & 3.584 & 3.878 \\
状態変化 & 失敗 & 33 & 1.758 & 3.818 \\
配置 & 成功 & 84 & 8.976 & 9.333 \\
配置 & 失敗 & 19 & 5.158 & 8.474 \\
\hline
\end{tabular}
\end{table}

失敗時に列が短くなるのは途中終了を含むためです．保存\texttt{attempts}は全18条件で保存列長と一致し，公開提案実装では最後の失敗試行，再生成，LLM要求は含まれません．参照長0は状態変化27件・配置2件で，長さ比の分母から除外して別集計しました．成功集合の違いにも注意し，手法間の共通成功タスクだけの比較を\path{Paper/analysis/common_success_comparisons.csv}に示します．本文の対応箇所は第5章の参照列長表と「評価指標」です．追加240件では保存長・実行試行数・失敗実行数・API要求数を分離して記録しました．

\subsubsection*{不明点B-4-2}
\begin{screen}
タスクの難易度を理解できる情報を提供してください。可能な範囲で、人間が同じ条件でタスクを行った時の成功率、アクション長などの情報を加えてください。完全なものでなくサンプリングした実験結果や、例示されたいくつかのタスクでのケーススタディでも構わないので、提案されたタスクの難易度を読者がある程度推定できるような情報が必要です。
\end{screen}
\subsubsection*{不明点B-4-2への回答}

ご指摘を受け，第5章「代表タスクのケース分析」に，状態変化2件・配置2件のケース分析を追加しました．本研究のタスクは，スイッチを点灯する，果物を冷蔵庫に入れるなど，人間にとって指示の意図や基本的な手順を日常知識から理解しやすいものです．一方，LLMを使ったタスクプランニングでは，日常的な手順の知識を，入力として与えられた個々の物体の状態や関係に対応付け，実行可能な行動列として具体化することが求められます．

例えば，すべてのスイッチを点灯する指示では，同名の物体をIDで識別し，点灯済みの物体と操作が必要な物体を区別して，未達の対象を順に処理する必要があります．初期状態ですべて点灯していれば，その時点で終了する判断も求められます．また，複数の果物を冷蔵庫に入れる指示では，対象への移動，把持，容器の開放，配置を，近接関係や保持状態に応じて順序付ける必要があります．LLMには，提示された環境知識からこれらの条件を読み取り，行動ごとに変化する状態を踏まえて，次の行動や終了を判断することが求められます．

追加したケース分析では，初期条件・ゴール・参照手順と，各段階でLLMに求められる判断を対応付けました．さらに，VH実行ログから，未達対象を残したまま終了する例や，再生成した行動が実行に失敗する例を示しました．人間にとっての難易度は日常的な手順理解という観点から説明し，LLMを使ったタスクプランニングにおける難しさを，具体的な判断過程と実際の失敗例によって示しました．

4ケースの説明内容について2026年9月6日に確認を受け，確認範囲を補足資料\texttt{Paper/author-cases.md}に記録しました．ケースの詳細と根拠ログは\texttt{Paper/vh/Cases.md}に示します．

\subsubsection*{不明点B-4-3}
\begin{screen}
提案手法のアプローチがどのような場合にうまく働き、どのような限界があったのか、を理解するための情報を提供してください。本論文では、提案手法全体の性能に関する統計値のみが記載され、議論されています。今回の実験では、タスクの自然言語記述から必要な状態や状況を正しく抜き出す自然言語処理的な部分に課題があるのか、プロンプトに与えた環境表現の情報量が推論に不足していたのか、アクションの粒度でうまく選択できていないのか、など、様々な側面での問題が想定されますが、実験結果で各ステップが実際にどのような入出力結果を出したのか、という情報がないため、提案手法の振る舞いが理解できません。こちらも、いくつかの例示でも構わないので、各ステップの出力例などを記載することが必要であると考えます。

もちろん、個々のコンポーネントごとの正解率、例えばprecondition判定の正解率など、それぞれのステップの性能が理解できる指標が追加されれば、読者に大きく役立つと考えます。
\end{screen}

\subsubsection*{不明点B-4-3への回答}

第5章「代表タスクのケース分析」に，未達対象が残った状態でEndを返す例と，前提条件を検証した後の再生成で不適切な操作が出力される例を追加しました．新しいVH試行では，各段階の完全グラフ，抽出知識，LLMへのメッセージ全文・生応答，生成候補，前提条件判定，違反説明，再生成，実行前後の状態とVH応答を保存しています．これにより，たとえばplumを保持している情報が終了判定の入力に存在したにもかかわらず誤終了した例を，情報抽出の欠落と区別して説明しました．

第5章「構成要素の比較」に，4タスク・240件の過程別集計を追加しました．C1のLLM判定を同じ候補・抽出辞書に対するJSON条件と照合した正答数はminiで34/43（79.1\%，誤受理4・誤棄却5），4oで45/54（83.3\%，誤受理9・誤棄却0）でした．情報不足は分母から除く規則ですが，このC1集計に該当例はありません．この集計の候補は軌跡ごとに異なります．

共通の固定問題を用いた別の静的評価では，前提条件の既知24入力を各3反復してminiは51/72，4oは53/72がJSONラベルと一致しました（各モデルで情報不足4入力の3反復，計12判定を除外）．これは構造化事実から作った入力に対する評価です．根拠は\path{Paper/nonvh/api_metrics.csv}です．

全10条件の途中時点を含む終了判定では，miniは496回中26回で未達なのにEndを返し，4oは600回中0回でした．達成済みなのに非Endは両モデル0回ですが，4oにはEnd/Continue以外の応答が2回あり，書式違反として別記しています．完全グラフ上のJSON条件がtrueだった実行でも，miniは376回中28回，4oは447回中37回でVHが失敗しました．これらはそれぞれ\path{Paper/vh_step4/termination_labels.csv}，\path{Paper/vh_step4/precondition_metrics.csv}，\path{Paper/vh_step4/execution_labels.csv}から確認できます．操作条件の判定，終了判定，VH実行制約を分けて扱う根拠としました．

同じ代表8候補について，候補実行直前の状態照合を実行条件とした評価も行いました．成功接頭列の再生後，記号一致・位置差最大0.001以下・姿勢差最大0.0001以下の場合だけ候補を実行し，2件が照合通過・VH成功でした．不一致6件には，人物のいないシーンへ候補時点の保存グラフを転送する方式を適用し，同じ照合でさらに2件が通過してVH成功しました．残る4件はFACING関係や物体の位置・姿勢・ON関係が一致せず候補未実行でした．両方式のAPI要求と通信エラーは0件でした．結果確認後，照合通過4候補の実行結果と未再現4件を報告しました．第5章の表「棄却候補8件の段階別再実行結果」に両方式の個別結果を示し，未再現範囲を本文で説明しました．

個別過程の例として，状態変化\texttt{test\_task6}のC1・GPT-4o・反復0・step 4では，候補\texttt{[SWITCHON] <lightswitch> (261)}にNoを返し，別の違反条件要求には\texttt{<lightswitch> (261) is OFF.}を返し，\texttt{[WALK] <lightswitch> (427)}へ再生成しています．ここでOFFという文は「満たされない条件」を列挙する要求への応答です．判定と違反条件の説明を分けた全文は，該当runの\path{events.jsonl}の\texttt{precondition}，\texttt{unmet}，\texttt{regeneration}に保存されています．

人手確認では，\texttt{precondition}応答を対象とした全8件「判断不能」（根拠「情報が不十分」）の票と，\texttt{unmet}の違反説明自体の確認を分けました．後者では要求・応答全文と対応状態を提示し，著者ugaiによる2026年9月7日の最終記入は「正しい」1件・「誤り」7件でした．評価者が当初は現在状態との一致で採点したと申告したため，別票で再確認し，AI支援で基準説明と個別例の整合性照会を複数回行いました．第5章「判定応答と違反説明の人手確認」に採点基準・確認手順・ラベル集計と記入例を示しました．著者記入の根拠全文は\path{Paper/human-unmet-review/semantic_review.csv}で公開しました．一部には未充足条件を列挙する要求と根拠の対応に曖昧さが残ります．判定応答の旧票も保持し，両評価対象を区別します．

代表候補の再実行は，LLM検証が棄却した59候補から8件を選び，元の成功接頭列を再生した後に候補を実行したものです．実際の選択規則は，タスク別にモデル・条件の組が異なる最初の候補を最大3件選ぶ方式で，状態変化\texttt{test\_task6}が3件，配置\texttt{test\_task1}が2件，配置\texttt{test\_task65}が3件です．初期充足の状態変化\texttt{test\_task1}には棄却候補がありません．選択対象はC1と検索変更条件です．

保存結果のVH成功は6/8件ですが，候補実行直前グラフと元の候補判定入力のグラフを照合すると，記号状態が一致したのは6/8件，さらに追加240件と同じ位置・姿勢の許容差を満たすのは2/8件でした．再生コードは候補時点の状態一致を実行前に検査していません．状態再現の制約を伴う追加再実行として報告します．選択表と結果は\path{Paper/vh_counterfactual/}，今回の事後照合は\path{Paper/reply-evidence/counterfactual_alignment.csv}に示します．元の240件とは別に集計しました．

\subsubsection*{不明点B-4-4}
\begin{screen}
比較対象ベースライン手法の実装について、情報を加えてください。
Huangらの方法は、Action TranslationやDynamic Exampleなどの工夫を加えていますが、これらがどのように本実験で実装されているのかが不明です。例えば、Dynamic Exampleは提案手法と同様の選択（タスク名による選択）でしょうか？再現できる程度の情報を記載してください。
\end{screen}
\subsubsection*{不明点B-4-4への回答}
第5章の比較条件・結果の説明と「評価間の条件差と記録範囲」に，BaselineはHuangらの著者公開実装（\url{https://github.com/huangwl18/language-planner}）を本研究の状態変化・配置タスクデータセットで実行して得た結果であると明記しました．手法の出典と数値の出所を区別し，実験結果の数値自体は変更していません．状態変化10/312成功，配置0/103成功ですが，状態変化の成功10件はすべて参照列長0で，非ゼロ長285件での成功は0件でした．

公開実装の出所と本研究のデータセットでの実行を確認し，現存する実験資料に基づいて説明範囲を明確にしました．実験時に使用したコミット・変更差分，GPT-4oへの接続処理，Action Translation・Dynamic Exampleの具体的なモデル・候補集合・設定，検索・停止・オブジェクト自動探索の詳細設定は，裏付けとなる追加資料を要する未確認項目として明示します．データ構築・分割，前提条件の原資料との対応，試行とコード・モデル・VH版の対応等についても，第5章「評価間の条件差と記録範囲」に確認済みの内容と資料上の制約を明記しました．実際に得た結果と実験資料の確認範囲を開示しました．

\subsubsection*{不明点B-4-5}
\begin{screen}
訓練データとテストデータがどの程度の類似タスクを含むのか、理解できる情報を記載してください。これは不足点2-2とも関係しますが、今回の実験においてはどのような設定を行ったのか、という情報が必要だと考えます。タスクの実行で、どの程度類似したタスクがプロンプトに埋め込まれたか、という例を記載してもらうのでも良いと思います。
\end{screen}
\subsubsection*{不明点B-4-5への回答}
第5章「データの分割方法」に検索類似度と対象・操作・配置先の一致率を，「実行設定と比較条件の詳細」に選択事例と検索規則を記載しました．検索類似度別の詳細集計は公開資料にまとめています．
第5章「データの分割方法」で，同名指示の非重複と構造の共有を区別しました．新しく再構成した検索について，テストインスタンスで重み付けした平均コサイン類似度は状態変化0.624，配置0.781です．配置103件では対象クラス一致率89.3\%，操作一致率54.4\%，配置先一致率10.7\%でした．これらは事例との構造的な一致率です．

例として``Put all plums in the fridge''には``Put all plums on the kitchentable''が類似度0.761で選択されます．対象は同じでも操作と配置先が異なり，入力環境に応じた変更が必要です．また``Turn off all tablelamps''にONの事例が類似度0.910で選ばれます．\path{Paper/analysis/retrieval_by_instance.csv}には事例ID・類似度・上位差・列長を保存しました．これは現在の再構成結果です．追加240件では保存済み選択を使った検索条件比較を行い，その結果を次の回答に示します．

\subsection*{コメントB-1}
\begin{screen}
Huangらの方法は、タスクなどをベクトル表現にして類似訓練例を得ています。本論文はその部分は「タスク名」で簡易化されているので、どの程度影響があるのか、その部分の性能を判断する指標があると良いと考えます。
\end{screen}
\subsection*{コメントB-1への回答}
検索方式が具体的に伝わるよう，タスク名のベクトル化と事例選択の手順を明記しました．実装は\texttt{sentence-transformers/all-MiniLM-L6-v2}でタスク名をベクトル化し，コサイン類似度最大の1名を選ぶ方式です．その名前のインスタンスから最長の参照列を選びます．検索入力はタスク名のみとし，シーン・初期状態・ゴールIDは検索対象外として整理しました．第4章の検索仕様と第5章「実行設定と比較条件の詳細」で明確化しました．

追加240件のC1を意味類似検索の基準として，random／fixed／noneの成功数を比較しました（表\ref{tab:reply_controls}）．miniでは意味類似9/12，random 8/12，fixed 7/12，none 3/12，4oではそれぞれ9/12，9/12，9/12，3/12でした．noneの成功3件はいずれも初期充足タスクです．ただし同じタスク・反復で対応させると，miniのrandomは意味類似に対して3件だけ成功へ変わり，4件だけ失敗へ変わっています．根拠は\path{Paper/vh_step4/paired_summary.csv}と\path{Paper/vh_step4/episodes.csv}です．randomは保存済み反復別選択を両モデルで共有し，fixedは固定1例，noneは事例メッセージを空にする条件でした．

\subsection*{コメントB-2}
\begin{screen}
p.3 左中央「プランニングを回帰的に行う」-$>$ 「再帰的」ではないでしょうか
\end{screen}

\subsection*{コメントB-2への回答}
ご指摘いただき，ありがとうございます．修正しました．

\subsection*{コメントB-3}
\begin{screen}
p.9 右「GPT-4o(...)の2種類のをLLMを用いて」 -$>$ 「を」不要
\end{screen}

\subsection*{コメントB-3への回答}
ご指摘いただき，ありがとうございます．修正しました．


\end{document}
